\subsection{熵变用初末态计算，熵产生还要看实际环境}
选一条可逆辅助路径，是为了计算状态函数的差，不是把真实过程宣称为可逆。图\ref{ch16:fig:s30-21-5}中的 $A\to B\to C$ 正是这种计算工具。系统熵变、热库熵变和二者的总熵变须分开；系统自己可以减熵，只要包括有关环境后的总熵不减。

\begin{origexamplebox}{同末态的两种膨胀，熵产生为何不同}\label{ch16:ex:teaching-two-paths}
$n$ 摩尔理想气体初态为 $(T_0,V_0)$。过程甲是在 $T_0$ 热库作用下可逆等温膨胀至 $2V_0$；过程乙是在隔热刚性容器中向真空自由膨胀至 $2V_0$。求气体、有关环境的熵变和总熵产生。忽略容器热容量、整体动能及势能的变化。
\end{origexamplebox}
\begin{xesolution}
乙过程既不传热也不向外做功，内能不变，理想气体末温仍为 $T_0$。所以甲、乙的气体初末态完全相同，均有
\[
 \Delta S_{\rm gas}=nc_v\ln\frac{T_0}{T_0}+nR\ln2=nR\ln2.
\]
甲从热库吸热 $Q=nRT_0\ln2$，热库熵变 $-Q/T_0=-nR\ln2$。接收功的理想机械装置不增加熵，故甲的总熵产生为零。
乙没有向环境传热，环境熵变为零，因此
\[
 S_{{\rm gen},\text{甲}}=0,\qquad
 S_{{\rm gen},\text{乙}}=nR\ln2>0.
\]
乙中 $Q=0$ 却有 $\Delta S_{\rm gas}>0$，并不违反熵定义，因为 $\Delta S=\int\delta Q_{\rm rev}/T$ 中的热量属于辅助可逆路径，而非真实自由膨胀的热量。若末体积趋于初体积，两者熵变都趋零；$nR$ 的单位为 $\mathrm{J/K}$，对数中的体积比无量纲。
\end{xesolution}

\subsection{有限热源不能始终当作恒温热库}
实际物体一边供热一边降温，热机每一小步可达到的效率也在改变。因此不能选一个平均温度代入卡诺效率再乘总热量。解决有限热源问题的顺序是：先明确装置是否循环复原，再写总能量和总熵条件，最后判断上界能否由可逆过程逼近。

\begin{origexamplebox}{两物体之间能够取出的最大功}\label{ch16:ex:teaching-finite-work}
两个物体的恒定热容量为 $C_1,C_2$，初温为 $T_h>T_c>0$。在它们之间运行热机，直到两者达到共同终温，热机最后恢复原态。整体与外界不交换热量，但可以输出机械功；忽略体积变化和其他能量改变。求最多能输出多少功，并与两物体直接接触比较。
\end{origexamplebox}
\begin{xesolution}
若共同终温为 $T_f$，热机不储存净能量，输出功为
\[
 W=C_1T_h+C_2T_c-(C_1+C_2)T_f.
\]
能量零点抵消，只涉及温度差。总熵不能减少：
\[
 C_1\ln\frac{T_f}{T_h}+C_2\ln\frac{T_f}{T_c}\ge0.
\]
把第一项的对数拆成 $\ln(T_f/T_c)-\ln(T_h/T_c)$，得到
\[
 (C_1+C_2)\ln\frac{T_f}{T_c}\ge C_1\ln\frac{T_h}{T_c},\qquad
 T_f\ge T_*=T_c\left(\frac{T_h}{T_c}\right)^{C_1/(C_1+C_2)}.
\]
所有对数和幂函数的自变量均无量纲。由于 $W$ 随 $T_f$ 增加而减小，最大功为
\[
 W_{\max}=C_1(T_h-T_*)+C_2(T_c-T_*).
\]
可在两物体的瞬时温度之间运行无穷小可逆循环，使总熵保持不变，直到温差趋零，所以该上界可以在可逆极限下逼近。图\ref{ch16:fig:s30-21-3}中的普通热机能流方向仍适用，但两端温度在不断变化。
直接接触时不输出功，终温是加权算术平均 $T_d=(C_1T_h+C_2T_c)/(C_1+C_2)\ge T_*$。若 $C_1=C_2=C$，则
\[
 T_*=\sqrt{T_hT_c},\qquad W_{\max}=C(\sqrt{T_h}-\sqrt{T_c})^2.
\]
两初温相等时最大功为零。直接传热没有损失总能量，却损失了温差所能提供的做功能力；不能把“能量守恒”误读为“有用功不变”。
\end{xesolution}

\begin{origexamplebox}{将有限物体冷却到环境温度以下所需的最少功}\label{ch16:ex:teaching-cooling-work}
热容量为常数 $C$ 的物体初温与环境热库同为 $T_0$。用循环复原的制冷机将物体冷却到 $T_1<T_0$，环境保持 $T_0$。求最少输入功。忽略体积变化及漏热，假定此温区无相变。
\end{origexamplebox}
\begin{xesolution}
在物体瞬时温度 $T$ 下再降低 $\dd T<0$，须取出热量 $\delta Q_c=-C\dd T>0$。可逆制冷系数为 $T/(T_0-T)$，所以最小元功为
\[
 \delta W_{\min}=(-C\dd T)\frac{T_0-T}{T}.
\]
系数随温度改变，必须积分：
\[
 W_{\min}=C\int_{T_1}^{T_0}\left(\frac{T_0}{T}-1\right)\dd T
 =C\left[T_0\ln\frac{T_0}{T_1}-(T_0-T_1)\right].
\]
物体总计放出 $Q_c=C(T_0-T_1)$，环境得到 $Q_h=Q_c+W_{\min}=CT_0\ln(T_0/T_1)$。因此物体熵变 $C\ln(T_1/T_0)$ 恰被环境的 $Q_h/T_0$ 抵消，验证了可逆极限。
令 $x=(T_0-T_1)/T_0$，小温差时 $W_{\min}\simeq CT_0x^2/2>0$，而不是与温差一次方成正比。真实有限速率制冷还会产生熵，所需功更大；本模型不能跨越热容变化或相变而直接外推至绝对零度。
\end{xesolution}
